%!TEX root = ../main.tex
%----------------------------------------------------------------------
%  Section 3: Structural Restrictions Catalog
%  Paper:   Theory-Informed Kernel Ridge Regression for Asset Pricing
%  Authors: Amedeo Andriollo and Juan F. Imbet
%----------------------------------------------------------------------

\section{Structural Restrictions Catalog}\label{sec:restrictions}

Each restriction in our framework originates from a specific economic model and imposes a distinct functional constraint on the mapping from observable characteristics to expected returns.
A consumption-based Euler equation with power utility implies a different penalty $C_j(f)$ than one with Epstein--Zin preferences, even though both restrict the same SDF moment condition, because the functional form of the stochastic discount factor $M_{j,t+1}$ differs.
Similarly, a production-based investment CAPM penalty differs from a $q$-theory penalty not merely in its parameter values but in which investment variables enter the mapping $g_j$ and how they interact.
The penalties are not interchangeable: each encodes a qualitatively different economic hypothesis about why expected returns vary across assets and over time.

We organize the fifty-six restrictions into eight families, grouped by economic mechanism: consumption-based Euler equations, production-based and investment restrictions, intermediary and institutional restrictions, information and learning restrictions, demand-system and market microstructure restrictions, variance and volatility restrictions, behavioral and sentiment restrictions, and macro-finance and term structure restrictions.
For each restriction, we specify three elements: the source model (the economic theory that implies it), the mathematical form of the restriction (the SDF, first-order condition, or equilibrium condition), and the penalty functional $C_j(f)$ through which it enters the Theory-KRR objective~\eqref{eq:tikrr} from Section~\ref{sec:framework}.
Because every penalty $C_j(f)$ depends on $f$ only through its evaluations at observed data points $\{f(\mathbf{X}_s)\}_{s=1}^n$, the representer theorem (Proposition~\ref{prop:representer}) applies to the full collection, and the estimator remains computationally tractable regardless of how many restrictions are active.


%======================================================================
\subsection{Consumption-Based Euler Equations}\label{sub:euler_restrictions}
%======================================================================

Consumption-based models derive asset-pricing restrictions from the representative agent's intertemporal optimality condition.
The Euler equation $\E[M_{t+1} R_{i,t+1}^{\mathrm{gross}} \mid \Fcal_t] = 1$ must hold for every asset $i$ and date $t$, where the SDF $M_{t+1}$ is determined by the agent's preferences and the consumption process.
Different preference specifications---power utility, recursive utility, habit formation, long-run risk---yield different SDFs and hence different penalties.
All thirteen restrictions in this family are Type~A (Euler equation) penalties as defined in Section~\ref{sub:euler}:
\begin{equation}\label{eq:euler_family}
C_j^{\mathrm{Euler}}(f) \;=\;
\sum_{t=1}^{T}
\left(
\frac{1}{N}\sum_{i=1}^{N}
M_j\bigl(\hat{\theta}_j, \mathbf{X}_t, \mathbf{X}_{t+1}\bigr)
\cdot
\bigl(f(\mathbf{X}_{i,t}) + R_t^f\bigr) - 1
\right)^2,
\end{equation}
where $\hat{\theta}_j$ denotes pre-estimated preference parameters obtained by GMM.
What distinguishes the thirteen restrictions is solely the form of $M_j$.

Table~\ref{tab:euler} lists the thirteen consumption-based restrictions. The thirteen consumption-based restrictions span several distinct economic channels.
Restrictions~1--5 vary preferences while maintaining i.i.d.\ or simple Markov consumption dynamics; what differs is how marginal utility depends on consumption growth and, in the habit models, on the history of consumption.
Power utility (Restriction~1) serves as the baseline, and each subsequent preference specification adds a mechanism---intertemporal substitution separated from risk aversion (Restriction~2), slow-moving external habit (Restriction~3), or direct dependence on lagged consumption (Restrictions~4--5).
These five SDFs share the same data requirements (aggregate consumption growth) but impose different functional relationships between $M_{t+1}$ and $\Delta c_{t+1}$.

Restrictions~6--9 modify the consumption \emph{process} rather than preferences alone.
Long-run risk (Restriction~6) introduces a persistent expected-growth component $x_t$ and stochastic volatility $\sigma_t$, both of which enter the SDF through Epstein--Zin recursion; the SDF depends on the news to $x_t$ and $\sigma_t^2$, not merely on realized consumption growth.
Adding jumps (Restriction~7) generates occasional large movements in volatility that are absent under Gaussian dynamics.
Rare disaster models (Restrictions~8--9) replace continuous volatility variation with low-probability, high-impact consumption contractions.
The economic content differs sharply: long-run risk explains the equity premium through intertemporal hedging demand, while disaster models explain it through the pricing of tail events.

Restrictions~10--13 extend the framework along further dimensions.
Precautionary savings (Restriction~10) introduces higher-order risk terms via the prudence coefficient, generating SDF dependence on conditional consumption variance.
Durable consumption (Restriction~11) adds a second consumption good with its own intratemporal elasticity.
Heterogeneous agents (Restriction~12) replaces the representative consumer with a continuum of agents whose idiosyncratic consumption risk enters the aggregate SDF through the cross-sectional variance of consumption growth.
The nonparametric Euler restriction (Restriction~13) abandons parametric preferences entirely, requiring only that some monotone decreasing function of consumption growth prices assets correctly---it serves as a benchmark that detects whether \emph{any} consumption-based SDF is empirically useful, without committing to a specific functional form.


%======================================================================
\subsection{Production-Based and Investment Restrictions}\label{sub:prod_restrictions}
%======================================================================

Production-based asset-pricing models derive expected returns from firms' intertemporal investment decisions.
The key insight is that firms equate the marginal cost of investment today to the expected discounted marginal benefit, and this first-order condition links expected stock returns to observable investment and profitability variables.
All ten restrictions in this family are Type~B (production FOC) penalties as defined in Section~\ref{sub:production}:
\begin{equation}\label{eq:prod_family}
C_j^{\mathrm{prod}}(f) \;=\;
\frac{1}{n}\sum_{s=1}^{n}
\bigl(f(\mathbf{X}_s) - g_j(\mathbf{Z}_s^{(j)})\bigr)^2,
\end{equation}
where $g_j$ maps firm-level investment characteristics $\mathbf{Z}_{i,t}^{(j)}$ to expected returns.
The functional form of $g_j$ and the variables that enter $\mathbf{Z}_{i,t}^{(j)}$ vary across restrictions.

Table~\ref{tab:production} lists the ten production-based restrictions. The production-based restrictions divide into three groups.
Restrictions~14--18 concern the technology of physical capital adjustment.
The investment CAPM (Restriction~14) and $q$-theory (Restriction~15) provide the broadest frameworks: expected returns are determined by firms' optimal investment behavior, with investment-to-capital ratio and profitability as the key state variables.
Restrictions~16--18 refine the adjustment cost specification.
Convex adjustment costs (Restriction~16) yield a smooth, monotone relationship between $I/K$ and expected returns.
Nonconvex adjustment costs (Restriction~17) and irreversibility (Restriction~18) break this smoothness: large adjustments incur fixed costs, and disinvestment is more expensive than investment.
These restrictions generate different functional forms for $g_j$---Restriction~16 predicts a linear or quadratic relationship in $I/K$, while Restrictions~17--18 predict threshold effects and asymmetries.

Restrictions~19--21 address financial frictions.
These models predict that expected returns reflect not only the marginal product of capital but also the shadow cost of financing constraints.
The three restrictions differ in which friction binds: collateral constraints (Restriction~19) link the premium to the ratio of capital to debt, equity issuance costs (Restriction~20) link it to leverage and internal cash flow, and costly external finance (Restriction~21) links it to the dividend payout and leverage.
All three predict that financially constrained firms earn higher expected returns, but the mapping from observables to the constraint's shadow value differs across models.

Restrictions~22--23 expand the set of productive inputs beyond physical capital.
Labor adjustment costs (Restriction~22) introduce the hiring rate as an additional state variable in the investment FOC, generating a two-dimensional mapping $(I/K, \mathrm{HN}) \mapsto \E[R]$.
Intangible capital (Restriction~23) adds R\&D spending and selling, general, and administrative expenses as measures of intangible investment, recognizing that knowledge capital and organizational capital are productive assets with their own adjustment costs.


%======================================================================
\subsection{Intermediary and Institutional Restrictions}\label{sub:intermediary_restrictions}
%======================================================================

Intermediary asset-pricing models posit that marginal investors are not households but financial intermediaries---broker-dealers, banks, hedge funds---whose pricing kernels reflect institutional constraints on leverage, capital, and funding.
When intermediaries are the marginal pricers, expected returns depend on intermediary-specific state variables: capital ratios, leverage, funding spreads, and balance-sheet capacity.
A linear factor model with factor $F_{j,t}$ implies the SDF $M_{j,t+1} = a_j + b_j F_{j,t+1}$, which is a special case of a linear SDF.
All eight restrictions in this family are therefore Type~A (Euler equation) penalties as defined in Section~\ref{sub:euler}:
\begin{equation}\label{eq:intermediary_family}
C_j^{\mathrm{Euler}}(f) \;=\;
\sum_{t=1}^{T}
\left(
\frac{1}{N}\sum_{i=1}^{N}
\bigl(\hat{a}_j + \hat{b}_j\, F_{j,t+1}\bigr)
\cdot
\bigl(f(\mathbf{X}_{i,t}) + R_t^f\bigr) - 1
\right)^2,
\end{equation}
where $F_{j,t}$ is the intermediary factor and $\hat{a}_j$, $\hat{b}_j$ are pre-estimated SDF parameters obtained by GMM.
This formulation requires only the factor time series---no pre-estimated betas are needed.

Table~\ref{tab:intermediary} lists the eight intermediary restrictions. The intermediary restrictions separate into models of intermediary health (Restrictions~24--25, 27) and models of market frictions that constrain intermediation (Restrictions~26, 28--31).
The intermediary capital ratio (Restriction~24) and leverage (Restriction~25) factors capture the state of broker-dealer balance sheets: when capital is depleted or leverage is high, the marginal value of wealth to the intermediary rises, increasing risk prices.
The dealer balance sheet factor (Restriction~27) constructs an equity capital ratio specifically for primary dealers, providing a more direct measure of the capacity constraint.
These three restrictions use different empirical proxies for the same conceptual channel---intermediary distress---and their relative $\mu_j$ values reveal which proxy captures the pricing-relevant variation.

Restrictions~26 and 28--31 focus on frictions in the trading environment.
Funding liquidity (Restriction~26) links expected returns to the cost and availability of short-term funding: when funding dries up, margins increase and leveraged positions are unwound, creating fire-sale dynamics.
Limits to arbitrage (Restriction~28) and slow-moving capital (Restriction~29) introduce a role for mispricing persistence---arbitrageurs face constraints (capital, risk limits, organizational frictions) that prevent immediate price correction.
The margin CAPM (Restriction~30) derives a two-beta model in which both market exposure and margin exposure are priced, while the leverage-constraint model (Restriction~31) predicts a compressed security market line because leverage-constrained investors overweight high-beta assets.
The key distinction: Restrictions~24--27 model \emph{who} the marginal pricer is, while Restrictions~28--31 model \emph{what constrains} the marginal pricer.


%======================================================================
\subsection{Information and Learning Restrictions}\label{sub:info_restrictions}
%======================================================================

Information and learning models generate expected-return variation from agents' incomplete knowledge about economic fundamentals.
Unlike consumption-based models, where the SDF is pinned down by preferences and the consumption process, these models derive pricing implications from the structure of beliefs---how agents learn, what they are uncertain about, and how they process information.
These restrictions generate Type~A penalties: Restrictions~32--35 specify a complete SDF from the learning model, while Restrictions~36--37 imply a linear SDF driven by a belief-related state variable (disagreement or information flow), which reduces to the Euler equation form $M_{j,t+1} = a_j + b_j F_{j,t+1}$.

Table~\ref{tab:information} lists the six information restrictions. The six information restrictions differ in \emph{what} agents are uncertain about and \emph{how} they process that uncertainty.
Bayesian learning (Restriction~32) generates time-varying risk premia because the posterior variance of growth declines as data accumulate, changing the conditional distribution of the SDF.
Regime learning (Restriction~33) produces a distinct mechanism: agents infer a latent state (expansion vs.\ recession), and uncertainty about the current regime amplifies return volatility and risk premia---the SDF depends on the belief $\pi_t$ rather than on a parameter estimate.
Parameter uncertainty (Restriction~34) prices the risk that the true data-generating process is unknown, generating a premium for model risk that shrinks only slowly with sample size.

Ambiguity aversion (Restriction~35) departs from Bayesian updating altogether: the agent entertains a set of plausible models $\mathcal{P}$ and evaluates each asset under the worst-case model.
The SDF is the likelihood ratio between the worst-case and reference models, generating pricing implications distinct from those of any single Bayesian model.
Disagreement (Restriction~36) and rational inattention (Restriction~37) shift focus from a single agent's beliefs to the cross-section of beliefs.
Disagreement models price the dispersion of beliefs as a risk factor (optimists hold long positions, pessimists short, and their trading generates a risk premium).
Rational inattention models generate endogenous information asymmetry: agents with limited processing capacity allocate attention optimally, and the resulting heterogeneity in information sets produces cross-sectional variation in expected returns.


%======================================================================
\subsection{Demand-System and Market Microstructure Restrictions}\label{sub:demand_restrictions}
%======================================================================

Demand-system models derive asset prices from the aggregation of heterogeneous investors' portfolio demands.
The equilibrium restriction is market clearing: aggregate demand for each asset must equal its supply.
This yields Type~C (demand-system equilibrium) penalties as defined in Section~\ref{sub:demand}:
\begin{equation}\label{eq:demand_family}
C_j^{\mathrm{demand}}(f) \;=\;
\sum_{t=1}^{T}
\left\|
\sum_{i} w_i^{(j)}\, D_i^{(j)}\bigl(\hat{\mathbf{f}}_t\bigr) - \mathbf{S}_t
\right\|^2,
\end{equation}
where $\hat{\mathbf{f}}_t = (f(\mathbf{X}_{1,t}), \ldots, f(\mathbf{X}_{N,t}))^\top$ and $\mathbf{S}_t$ denotes market-capitalization shares.
Restrictions~41--43 involve microstructure frictions and take modified forms as described below.

Table~\ref{tab:demand} lists the six demand-system and microstructure restrictions. The demand-system restrictions (38--40) share the market-clearing equilibrium condition but differ in the structure of investors' demand functions.
Restriction~38 models demand as a function of observable characteristics and a latent demand shock, estimating the demand system directly from portfolio holdings data.
Restriction~39 derives demand from investors' optimality conditions with characteristic-based demand elasticities, providing a micro-founded demand function.
Restriction~40 focuses on the \emph{aggregate} demand elasticity: when total market demand is inelastic, flow-driven trading has outsized price impact, and expected returns reflect the accommodation of demand shocks rather than fundamental risk.

Restrictions~41--43 address frictions that prevent instantaneous and costless trading.
Search frictions (Restriction~41) apply primarily to over-the-counter markets, where finding a counterparty takes time; the expected-return penalty includes the search intensity and the buyer's bargaining power.
Fire sale externalities (Restriction~42) generate a Type~A Euler penalty with $M_{42,t+1} = \hat{a}_{42} + \hat{b}_{42}\,\mathrm{FSP}_{t+1}$: when multiple leveraged institutions simultaneously liquidate positions, the fire-sale pressure factor enters the SDF.
Short-sale constraints (Restriction~43) restrict the adjustment mechanism on the negative side: when pessimists cannot short, asset prices reflect only optimists' valuations, and the degree of overpricing depends on the divergence of opinions and the binding intensity of the constraint.


%======================================================================
\subsection{Variance and Volatility Restrictions}\label{sub:option_restrictions}
%======================================================================

Variance and volatility models link expected returns to measures of conditional variance and the compensation investors demand for bearing variance risk.
These two restrictions use publicly available data---the CBOE Volatility Index (VIX) as a measure of implied variance and daily CRSP returns to compute realized variance---rather than individual equity option data, making them applicable to the broad cross-section over the post-1990 period.
Both restrictions generate Type~B penalties that map volatility-related state variables to expected returns.

Table~\ref{tab:options} lists the two variance and volatility restrictions. The two variance and volatility restrictions capture complementary aspects of volatility pricing.
The variance risk premium (Restriction~44) uses the spread between implied variance (from VIX) and expected realized variance (computed from daily CRSP returns within each month) as a predictor of aggregate returns; the penalty penalizes discrepancies between $f(\mathbf{X}_t)$ and the VRP-implied expected return.
\citet{bollerslevtauchenzhou2009} show that VRP is a strong predictor of market returns at horizons from one to six months, and the restriction encodes this predictive relationship.
Stochastic volatility (Restriction~45) uses the conditional variance $V_t$ from a Heston-type model, where VIX serves as a proxy for the latent conditional variance and the variance process carries its own risk premium distinct from the equity premium.
Together, these restrictions ensure that the estimator respects the well-documented relationship between variance dynamics and expected returns without requiring granular option-level data.


%======================================================================
\subsection{Behavioral and Sentiment Restrictions}\label{sub:behavioral_restrictions}
%======================================================================

Behavioral models derive expected-return patterns from departures from rational expectations or expected utility.
These restrictions do not derive an SDF from a representative agent's Euler equation; instead, they predict cross-sectional or time-series return patterns as functions of behavioral variables---past returns, capital gains, sentiment indices.
Most generate Type~B penalties, where $g_j$ maps behavioral state variables to expected returns.

Table~\ref{tab:behavioral} lists the six behavioral restrictions. The behavioral restrictions exploit distinct psychological mechanisms.
Prospect theory (Restriction~46) modifies the SDF with a kink at the reference point: losses are weighted more heavily than gains by a factor $\lambda > 1$, generating a first-order risk premium for assets whose returns covary with movements around the reference point.
This is the only behavioral restriction that directly specifies an SDF, placing it closest to the consumption-based framework.

Restrictions~47--49 generate cross-sectional return predictability from belief formation biases.
Overconfidence (Restriction~47) produces momentum at short horizons (agents overreact to private signals, pushing prices beyond fundamentals) and reversal at long horizons (prices revert as public information arrives).
The disposition effect (Restriction~48) predicts return continuation among stocks with large unrealized gains and return reversal among stocks with large unrealized losses, because investors' reluctance to realize losses impedes price adjustment.
Extrapolative expectations (Restriction~49) model agents who weight recent returns too heavily when forming expectations, generating a specific functional form for the relationship between past returns and expected future returns.
These three restrictions share momentum-related variables as inputs but predict different functional relationships between past returns and expected future returns.

Investor sentiment (Restriction~50) operates at the aggregate level: the Baker--Wurgler sentiment index predicts that, conditional on high sentiment, speculative stocks (small, unprofitable, high-volatility) earn low subsequent returns, while safe stocks are largely unaffected.
This generates a Type~A Euler equation penalty where $M_{50,t+1} = \hat{a}_{50} + \hat{b}_{50}\,\mathrm{Sent}_{t+1}$ is a linear SDF in the sentiment factor.
Salience theory (Restriction~51) predicts mispricing at the individual asset level: stocks with salient (extreme) upside payoffs are overpriced because agents overweight the probability of salient outcomes, generating a negative relationship between salience and expected returns.


%======================================================================
\subsection{Macro-Finance and Term Structure Restrictions}\label{sub:macro_restrictions}
%======================================================================

Macro-finance models link expected returns on equities and bonds to macroeconomic state variables: the yield curve, inflation expectations, monetary policy, and consumption-wealth ratios.
These restrictions enter as Type~A or Type~B penalties depending on whether the model implies a linear SDF driven by a macro factor (Type~A) or a direct functional mapping from macro variables to expected returns (Type~B).

Table~\ref{tab:macro} lists the five macro-finance restrictions. The macro-finance restrictions link equity expected returns to macroeconomic and fixed-income state variables through five distinct channels.
The affine term structure restriction (Restriction~52) imposes no-arbitrage across maturities, using the yield curve state vector $\mathbf{y}_t$ (level, slope, curvature, and potentially macro factors) to constrain the time-varying equity risk premium.
This restriction is structurally different from the others because it derives from the bond market and imposes cross-market no-arbitrage.

Restrictions~53--54 transmit monetary and inflation shocks to equity prices.
The inflation risk premium (Restriction~53) uses the spread between nominal and real yields to extract the compensation investors demand for bearing inflation uncertainty; this premium predicts equity returns because inflation uncertainty covaries with real economic risk.
The monetary policy restriction (Restriction~54) identifies the surprise component of federal funds rate decisions (using fed funds futures) and maps it to equity returns, exploiting the direct channel from policy rates to discount rates.

Restrictions~55 and~56 use slow-moving macroeconomic aggregates.
The consumption-wealth ratio $cay$ (Restriction~55) exploits cointegration among consumption, asset wealth, and labor income: when consumption is high relative to wealth, expected returns are high because the ratio mean-reverts through subsequent returns.
The excess bond premium (Restriction~56) isolates the component of credit spreads that reflects risk appetite and financial intermediary health, net of default risk; it predicts equity returns because deteriorating credit conditions signal tightening financial constraints economy-wide.


%======================================================================
\subsection{Summary and Discussion}\label{sub:restriction_summary}
%======================================================================

The catalog comprises fifty-six restrictions organized into eight families.
Each restriction imposes a distinct functional form on the mapping from observable characteristics $\mathbf{X}_{i,t}$ to expected returns $f(\mathbf{X}_{i,t})$, and each translates into a specific penalty functional $C_j(f)$ within the Theory-KRR objective~\eqref{eq:tikrr}.
The penalties are heterogeneous in structure: Type~A penalties enforce Euler equation moment conditions and require specifying the SDF---either from a structural model or from a linear factor representation $M_{t+1} = a + b F_{t+1}$ (Restrictions~1--13, 24--37, 42, 46, 50, 53, 56); Type~B penalties equate expected returns to functions of investment, behavioral, or macro variables through first-order conditions (Restrictions~14--23, 41, 43, 44--45, 47--49, 51--52, 54--55); and Type~C penalties enforce market-clearing equilibrium conditions from demand-system models (Restrictions~38--40).
Despite this heterogeneity, the framework in Section~\ref{sec:framework} accommodates all fifty-six penalties within a single objective: each $C_j(f)$ depends on $f$ only through its evaluations at observed data points, the representer theorem applies, and cross-validation selects the vector $\bm{\mu}$ that balances predictive fit against theoretical guidance.

A property of the catalog that makes the $\mu_j$ ranking meaningful is that many restrictions share observable input variables but impose different functional relationships.
Consumption growth enters Restrictions~1--13 through thirteen different SDFs.
Investment-to-capital ratio enters Restrictions~14--18 through five different adjustment cost specifications.
Past returns enter Restrictions~47--49 through three different behavioral mechanisms.
When two restrictions share inputs but one receives a higher $\mu_j$ in cross-validation, the data have identified which functional relationship between those inputs and expected returns is more empirically useful.
This is the sense in which the cross-validated multipliers $\bm{\mu}$ produce a ranking of economic theories: they rank not just which state variables matter, but which \emph{economic mechanisms linking state variables to expected returns} are supported by the data.
